% \newpage
\section*{Appendix}
\label{sec: Appendix}

% %%%%%%%%%%%%%%%%%%%%%%%%%%%%%%%%%%%%%%%%%%%%%%%%%%%%%%%%%%%%%%
% %%%%%%%%%%%%%%%%%%%%%%%%%%%%%%%%%%%%%%%%%%%%%%%%%%%%%%%%%%%%%%
% \subsection{Bellman operators}

% \textcolor{red}{add definition in appendix}
% The inherent recursive structure of action-value functions can be elegantly formalized using the Bellman evaluation operator $T_\pi$ and the Bellman optimality operator $T$. 
% % acting on the space of action-value functions $\qset$.
% We define these operators as
% \begin{align}
%     (T_\pol q) (s,a) &\coloneqq \reward(s,a) + \discount \sum_{s' \in \sset} p(s' \mid s,a) q(s', \pol(s')) ,
%     \\
%     (T q)(s,a) &\coloneqq \reward(s,a) + \discount \sum_{s' \in \sset} p(s' \mid s,a) \max_{a' \in \aset(s')} q(s',a') .
% \end{align}
% We can concisely write $T_\pi q = r + \discount P_\pi q$,
% % \begin{align*}
% %     T_\pi q = r + \discount P_\pi q ,
% % \end{align*}
% where $P_\pi$ is the transition operator defined as
% \begin{align*}
%     (P_\pi q)(s,a) \coloneqq \sum_{s' \in \sset} p(s' \mid s,a) q(s', \pol(s')) .
% \end{align*}
% For finite MDPs with bounded rewards, the operators $T_\pi$ and $T$ exhibit two fundamental properties.
% They are monotone
% and they are $\gamma$-contractions in the $L_\infty$-norm
% provided the discount factor satisfies $\gamma \in [0,1)$
% or the MDP is episodic (i.e. every policy eventually reaches a terminal state).
% In other words, for every policy $\pi$ and any two action-value functions $q_1$ and $q_2$,
% \[
% q_1 \leq q_2 
% \quad\implies\quad
% T_\pi q_1 \leq T_\pi q_2 
% \]
% and
% \[
% \| T_\pi q_1 - T_\pi q_2 \|_\infty \leq \gamma \|q_1 - q_2 \|_\infty,
% \]
% and similarly for the optimality operator $T$.
% These properties guarantee that $q_\pi$ and $\qopt$ are the \textit{unique} fixed points of $T_\pi$ and $T$, respectively~\citep{Puterman1994MarkovProgramming.}.
% As a result, a policy $\pi$ is optimal if and only if it satisfies the Bellman optimality equation
% \begin{align}
% \label{eq: bellman optimality equation}
%     q_\pi = T q_\pi .
% \end{align}
% or, given $q_\pi = T_\pi q_\pi$, if it is greedy with respect to its own action values, i.e. $T_\pi q_\pi = T q_\pi$.


%%%%%%%%%%%%%%%%%%%%%%%%%%%%%%%%%%%%%%%%%%%%%%%%%%%%%%%%%%%%%%
%%%%%%%%%%%%%%%%%%%%%%%%%%%%%%%%%%%%%%%%%%%%%%%%%%%%%%%%%%%%%%
% \subsection{Proofs}

%%%%%%%%%%%%%%%%%%%%%%%%%%%%%%%%%%%%%%%%%%%%%%%%%%%%%%%%%%%%%%
\paragraph{Proof of \autoref{thm: properties of noise}}

% \begin{proof}
Fix a time $k\ge 0$, a state $s \in \mathcal S$ and an action $a \in \mathcal A(s)$.
Let $\mathcal F_k'$ represent the available information right after selecting the policy $\pi_k$ but before sampling the initial state-action pair, so $\mathcal F_k \subseteq \mathcal F_k'$.

By definition,
\begin{align}
\label{eq: noise def in appendix proof}
w_k(s,a) = u_k(s,a) v_k(s,a) + \big( u_k(s,a)-\mu_k(s,a) \big) \big( q_{\pi_k}(s,a) - q_k(s,a) \big).
\end{align}
Since episode returns are unbiased under initial- or first-visit updates, 
and $u_k(s,a) v_k(s,a) = 0$ on $\{u_k(s,a)=0\}$,
it follows that
\begin{align}
\label{eq: unbiased noise p1}
\mathbb E[ u_k(s,a) v_k(s,a) \mid \mathcal F_k']
=
\mathbb E \big[ 
    \mathbf 1_{\{u_k(s,a)=1\}} \mathbb E[ v_{k}(s,a) \mid \mathcal F_k', u_k(s,a)=1] 
    % +
    % \mathbf 1_{\{u_k(s,a)=0\}} \mathbb E[ v_{k}(s,a) \mid \mathcal F_k', u_k(s,a)=o] 
    \mid \mathcal F_k' \big]
= 0 .
\end{align}
Also, $\mathbb E[u_k(s,a) \mid \mathcal F_k'] = \mu_k(s,a)$, and $\mu_k$ and $q_{\pi_k}-q_k$ are $\mathcal F_k'$-measurable. Hence,
\begin{multline}
\label{eq: unbiased noise p2}
\mathbb E \!\left[
\bigl( u_k(s,a)-\mu_k(s,a) \bigr) \bigl( q_{\pi_k}(s,a)-q_k(s,a) \bigr)
\mid \mathcal F_k'
\right]
\\
=
\big( q_{\pi_k}(s,a)-q_k(s,a) \big) 
\big( \mathbb E\!\left[ u_k(s,a) \mid \mathcal F_k'
\right]
-\mu_k(s,a) \big)
=
0 .
\end{multline}
Summing \eqref{eq: unbiased noise p1} and \eqref{eq: unbiased noise p2} yields $\mathbb E[w_k(s,a) \mid \mathcal F_k'] = 0$.
Applying the tower property then proves \eqref{eq: muw_mds}:
\[
\mathbb E [ w_{k}(s,a) \mid \mathcal F_k ]
= \mathbb E \!\left[ \mathbb E [ w_{k}(s,a) \mid \mathcal F_k'] \mid \mathcal F_k \right]
= 0 .
\]

Finally, since $u_k^2(s,a) \le 1$ and $(u_k(s,a) - \mu_k(s,a))^2 \le 1$,
using $(x+y)^2 \le 2 x^2 + 2 y^2$ on \eqref{eq: noise def in appendix proof} yields
\begin{align*}
w_k^2(s,a) 
&\le 2 v_k^2(s,a) + 2 (q_{\pi_k}(s,a)-q_k(s,a) )^2 
\\
&\le 2 v_k^2(s,a) + 2 ( 2 q_{\pi_k}^2(s,a) + 2 q_k^2(s,a) ) .
\end{align*}
Since all action-values are bounded,
there exists a constant $c_\polset < \infty$ such that $\|q_\pi\|^2 \leq c_\polset$ for every policy $\pi \in \polset$. As a result,
\begin{align*}
\mathbb E \!\big[ w_{k}^2(s,a) \mid \mathcal F_k \big]
&\le
2 c_v + 2 \, \mathbb E \!\left[ ( q_{\pi_k}(s,a)-q_k(s,a) )^2 \mid \mathcal F_k \right]
\\
&\le
2 c_v + 4 c_\polset + 4 \|q_k\|^2
\end{align*}
So \eqref{eq: muw_second_moment} holds for example with $c_w := \max \{ 2 c_v + 4 c_\polset, 4\}$.
% \end{proof}
\hfill \BlackBox


%%%%%%%%%%%%%%%%%%%%%%%%%%%%%%%%%%%%%%%%%%%%%%%%%%%%%%%%%%%%%%
\paragraph{Proof of \autoref{thm: bounded limit sets}}

% \begin{proof}
Let $(q_k)_{k \ge 0}$ and $(\pol_k)_{k \ge 0}$ be solutions of the difference inclusion in \eqref{eq: mcopi sa} with initial condition $q_0$, learning rates $(\lrate_k)_{k \ge 0}$, update distributions $(\update_k)_{k \ge 0}$ and stochastic perturbations $(w_k)_{k \ge 0}$.
    
    By assumption, there exists a finite time $K$ such that
    the sequence $(\pol_k)_{k \ge K}$ only contains policies in $\polset_\infty$.
    % the estimate $q_k$ is updated towards action values $q_{\pi_k}$ where $\pi_k \in \polset_\infty$.
    Thus for all $k \ge K$, each state-action pair of the estimate $q_k$ is updated towards a value $q_{\pol_k}(s,a)$ that satisfies
    \begin{align}
        \min_{\pi \in \polset_\infty} q_\pi(s,a) \leq q_{\pi_k}(s,a) \leq \max_{\pi \in \polset_\infty} q_\pi(s,a) .
    \end{align}
    Therefore, define the sequences $(x_k)_{k \ge K}$ and $(y_k)_{k \ge K}$ as
    \begin{align}
        x_\kpo &= x_k + \lrate_k \big( \update_k(s,a) 
        ( \min_{\pi \in \polset_\infty} q_\pi(s,a) - x_k ) + w_k(s,a) \big) ,
        \\
        y_\kpo &= y_k + \lrate_k \big( \update_k(s,a) ( \max_{\pi \in \polset_\infty} q_\pi(s,a) - y_k ) + w_k(s,a) \big) ,
    \end{align}
    and initialised at $x_K = y_K = q_K(s,a)$. From time $K$ onwards we have
    \begin{align*}
        x_k \leq q_k(s,a) \leq y_k .
    \end{align*}
    Proposition~4.1 and Example~4.3 in \citep{Bertsekas1996Neuro-DynamicProgramming} show that $x_k$ and $y_k$ converge almost surely to $\min_{\pi \in \polset_\infty} q_\pi(s,a)$ and $\max_{\pi \in \polset_\infty} q_\pi(s,a)$, respectively.
    % because the noise satisfies conditions \eqref{eq: muw_mds} and \eqref{eq: muw_second_moment}.
    As a result, we conclude that, almost surely,
    \begin{align*}
        &\liminf_{k \to \infty} q_k(s,a)
        \ge \min_{\pi \in \polset_\infty} q_\pi(s,a)
        % \leq q_k(s,a) 
        \\
        &\limsup_{k \to \infty} q_k(s,a)
        \le \max_{\pi \in \polset_\infty} q_\pi(s,a)
    \end{align*}
    for every state-action pair,
    which proves the first part of the lemma.

    Further,
    given that there are a finite number of deterministic policies, $q_k$ converges almost surely to a bounded set.
    As a result, 
    % that sequence is almost surely bounded, meaning that 
    its norm over all $k \geq 0$ is almost surely finite. So there exists an almost surely finite constant $c_Q$ such that $\sup_{k \geq 0} \|q_k\|_\infty^2 < c_Q$.
% \end{proof}
\hfill \BlackBox


%%%%%%%%%%%%%%%%%%%%%%%%%%%%%%%%%%%%%%%%%%%%%%%%%%%%%%%%%%%%%%
% \subsection{Proof of \autoref{thm: effects of noise decrease}}

% \textcolor{red}{need to fix a state-action pair, otherwise x is not scalar}

% % \begin{proof}[Proof of \autoref{thm: effects of noise decrease}]
% Let $(q_k)_{k \ge 0}$ be a solution of the difference inclusion in \eqref{eq: complete difference inclusion} with initial condition $q_0$, learning rates $(\lrate_k)_{k \ge 0}$, update distributions $(\update_k)_{k \ge 0}$ and stochastic perturbations $(w_k)_{k \ge 0}$.
    
%     % \autoref{eq: properties of noise} shows that there exist two positive constants $c_1$ and $c_2$ such that, for every state-action pair and $k \ge 0$,
%     % \begin{align*}
%     %     &\E \big[ \update_k(s,a) w_k(s,a) \big] = 0 ,
%     %     \\
%     %     &\E \big[ \update_k^2(s,a) w_k^2(s,a) \big] \leq c_1 + c_2 \|q_k\|^2 ,
%     % \end{align*}
%     % while \autoref{thm: bounded solutions} establishes that there exists almost surely a positive constant $c_3$ such that, for all $k \geq 0$,
%     % \begin{align*}
%     %     \|q_k\|^2 < c_3 .
%     % \end{align*}
    
%     Define a sequence of partial sums $(x_k)_{k \ge 0}$ as 
%     \begin{align}
%         x_k = \sum_{i = 0}^{k} \lrate_{i} w_{i} .
%         % x_\kpo = x_k + \lrate_k \update_k w_k
%     \end{align}
%     % Let $\mathcal{F}_k$ be the filtration defined in \eqref{eq: filtration_def}. 
%     By \autoref{thm: properties of noise}, the perturbations are conditionally unbiased, meaning that $\mathbb{E}[\lrate_k w_k \mid \mathcal{F}_k] = 0$. Therefore, $(x_k)_{k \ge 0}$ is a martingale adapted to the filtrations $(\mathcal{F}_{k})_{k \ge 0}$.
    
%     \autoref{thm: properties of noise} and \autoref{thm: bounded limit sets} 
%     also imply that there exist almost surely finite constants $c_w, c_Q \in [0, \infty)$ such that the variance of the martingale increments satisfy
%     \begin{align*}
%     \E \left[(x_k - x_{k-1})^2 \mid \mathcal{F}_k \right] 
%     =
%     \alpha_k^2 \E [w_k^2 \mid \mathcal{F}_k] 
%     % \le \lrate_k^2 C(1 + \|q_k\|^2)
%     \le \lrate_k^2 c_w(1 + c_Q) .
%     \end{align*}
%     Thus if the learning rates satisfy $\sum_{k=0}^\infty \alpha_k^2 < \infty$, the sum of the variances is finite:
%     \[
%     \sum_{k=1}^\infty \mathbb{E}[(x_k - x_{k-1})^2] \le c_w(1 + c_Q) \sum_{k=1}^\infty \alpha_k^2 < \infty .
%     \]
%     Altogether, $(x_k)_{k \ge 0}$ is a martingale with bounded $L^2$ variations. Thus Doob's Martingale Convergence Theorem indicates that $x_k$ converges almost surely to a finite limit. Furthermore, since the sequence of partial sums converges, the individual terms $\alpha_k w_k$ almost surely converge to zero.
% % \end{proof}

%%%%%%%%%%%%%%%%%%%%%%%%%%%%%%%%%%%%%%%%%%%%%%%%%%%%%%%%%%%%%%
% \subsection{Proof of \autoref{thm: steps become small}}


% \begin{proof}[Proof of \autoref{thm: steps become small}]
% The Robbins-Monro conditions and the boundedness of the update distributions ensure that the sequence $(\lrate_k \update_k)$ converges to zero for all state-action pairs.
% Additionally, \autoref{thm: bounded limit sets} implies that solutions of MC-O-PI converge almost surely to the region
% \begin{align}
%     \big\{q \in \qset : q_{\worst{\pi}} \le q \le \qopt \big\} .
% \end{align}
% where $q_{\worst{\pi}}$ and $\qopt$ are the action values of the worst and the best policies, respectively.
% Therefore, for every $\epsilon > 0$, there exists almost surely a time step $K'$ such that, for all $k \geq K'$,
% \begin{align*}
%  \|\lrate_k \update_k (q_{\pol_k} - q_k)\| 
%  &\leq \|\lrate_k \update_k \| \|q_{\pol_k} - q_k\| 
%  \\
%  &\leq \|\lrate_k \update_k \| \|\qopt - q_{\worst{\pol}}\|
%  \\
%  &\leq \frac{\epsilon}{2} .
% \end{align*}
% Further, \autoref{thm: effects of noise decrease} demonstrates that there exists almost surely a time step $K''$ such that, for all $k \geq K''$,
% \begin{align*}
%  \|\lrate_k \update_k w_k \| \leq \frac{\epsilon}{2} .
% \end{align*}
% As a result, for every $\epsilon > 0$,
% there exists almost surely a time step $K \geq \max\{K', K''\}$ such that, for all $k \geq K$,
% \begin{align*}
%     \| q_\kpo - q_k \| 
%     &= \| \lrate_k \update_k (q_{\pol_k} - q_k + w_k) \|
%     \\
%     &\leq \| \lrate_k \update_k (q_{\pol_k} - q_k) \| + \| \lrate_k \update_k w_k \|
%     \\
%     &\leq \epsilon ,
% \end{align*}
% indicating that the sequence $(\|q_\kpo - q_k\|)$ converges almost surely to zero.
% \end{proof}


%%%%%%%%%%%%%%%%%%%%%%%%%%%%%%%%%%%%%%%%%%%%%%%%%%%%%%%%%%%%%%
\paragraph{Proof of \autoref{thm: convergence deterministic mcopi}}

% \begin{proof}
Let $(\pi_k)_{k \ge 0}$ be the sequence of policies associated with a solution of \eqref{eq: mcopi sa}
and let $(t_n)_{n \ge 0}$ be the corresponding transition times.
By assumption, the update distributions are uniform over all actions in each state and satisfy exploring starts.
\autoref{thm: improving policies} thus ensures that the sequence of action values $(q_{\pi_k})_{k \ge 0}$ is componentwise non-decreasing.

Since the policy set is finite, the subsequence $(q_{\pi_{t_n}})_{n \ge 0}$ must settle after finitely many transitions.
In other words, there exists a finite index \(N \ge 0\) such that
$t_N < \infty$ and $t_{N+1} = \infty$.
Moreover, policy $\pi_{t_N}$ must be optimal. Otherwise the modified Robbins-Monro and exploring starts conditions would force $q_k$ to converge to $q_{\pi_{t_N}}$, which in turn would cause the policy to improve and thus $t_{N+1}<\infty$ (as in \autoref{thm: suboptimal blocks finite}).

As a result, for every $k \ge t_N$, the solution of \eqref{eq: deterministic mcopi} satisfies
\[
q_{k+1} = q_k + \alpha_k \mu(\sigma_k, \best{\pi}, f_k)( \qopt - q_k) .
\]
The Robbins-Monro and strict exploring starts conditions then guarantee that $q_k$ converges to $\qopt$ \citep[Proposition~4.1]{Bertsekas1996Neuro-DynamicProgramming}.
% \end{proof}
\hfill \BlackBox
